\section{BHSM interpretation of the reconstructed exterior}
\label{sec:bhsm-interpretation}

The ordinary lens fit is part of the BHSM interpretation used here.
It reconstructs the effective gravitational exterior of the compact
configuration: its monopole, location, translation and admitted projected
structure. An observable departure from GR is not required for that
interpretation. The recovered BHSM source records explicitly admit an
Einstein/GHY compact-source sector, in which boundary and conventional
gravitational descriptions represent the same exterior contribution.
We use that declared sector without promoting it into a unique
microscopic derivation of every fitted parameter.

Let $\mathcal C$ denote a configuration described in the BHSM framework,
and let $W_{\mathcal C}=\Phi_{\mathcal C}+\Psi_{\mathcal C}$ be its
Weyl potential in the stated weak-metric convention. The forward map is
\begin{equation}
 \mathcal C\longmapsto W_{\mathcal C}
 \longmapsto\psi_{\mathcal C}
 \longmapsto(\boldsymbol\theta_+,\boldsymbol\theta_-,
                  \ln R,\Delta t),
 \label{eq:bhsm-optical-interpretation}
\end{equation}
with the supplied distances, clock, angular frame and source positions.
When a source convention writes
$W_{\rm topo}=(\Phi_{\rm topo}+\Psi_{\rm topo})/2$, its contribution to
our full Weyl potential is $2W_{\rm topo}$; this conversion creates no
second physical response.
The fit constrains the projected response represented by
$\psi_{\mathcal C}$. Two descriptions giving the same projected potential
on the relevant image neighborhoods give the same positions,
magnification determinants and relative Fermat delays. Their likelihood
and fitted track are consequently identical. This is a physical
interpretation of the reconstructed exterior, rather than a demand for
an additional force or a second mass contribution.

\begin{table}[tbp]
\centering\small
\caption{Interpretation dictionary for the ordinary optical fit. All
entries remain conditional on the supplied geometry and finite lens model.}
\label{tab:bhsm-interpretation}
\begin{tabular}{@{}p{.23\linewidth}p{.71\linewidth}@{}}
\toprule
Fitted quantity & BHSM interpretation and what is constrained\\\midrule
Optical mass $M_{\rm opt}$ & Effective gravitational monopole of the
represented configuration. It is the point-exterior lens coefficient,
not a measurement of a microscopic constituent mass or a surface energy
without its matching relation.\\
Reference center $\boldsymbol c_0$ & Location of that optical
configuration in the supplied tangent-plane frame.\\
Observer proper motion $\boldsymbol\mu_{\rm obs}$ & Translation of
that same configuration per observer year. It does not measure internal
envelope evolution or select a causal relaxation law.\\
Compact quadrupole & Projected nonspherical structure of the same
exterior; an envelope may be represented by this term when it produces
that potential. The coefficient does not uniquely identify its cause.\\
Shear and convergence & The admitted projected environmental response.
Their interpretation within a BHSM environment does not assign them
uniquely to the tracked compact object.\\\bottomrule
\end{tabular}
\end{table}

The spherical boundary interpretation can be made explicit without
changing the lens equations. In the static Schwarzschild member of the
declared Einstein sector, flat-reference Brown--York energy obeys
\begin{equation}
 E_{\rm BY}(r)=\frac{c^4r}{G}
       \left(1-\sqrt{1-\frac{2GM_{\rm opt}}{c^2r}}\right),\qquad
 M_{\rm opt}=\frac{E_{\rm BY}}{c^2}
       -\frac{G E_{\rm BY}^2}{2c^6r},\quad
 r>\frac{2GM_{\rm opt}}{c^2}.
 \label{eq:bhsm-boundary-mass-interpretation}
\end{equation}
This is the established gravitational boundary construction
\cite{BrownYork1993}, also recovered in the BHSM source record.
The corresponding principal spatial-curvature deficit relative to a flat
sphere and its inverse mass relation are
\begin{equation}
 \delta k_g(r)=\frac{1-\sqrt{1-2GM_{\rm opt}/(c^2r)}}r,\qquad
 M_{\rm opt}=\frac{c^2}{G}
       \left(r^2\delta k_g-\frac12r^3\delta k_g^2\right).
 \label{eq:bhsm-curvature-mass-interpretation}
\end{equation}
For fixed declared $r$, let $f=1-2GM_{\rm opt}/(c^2r)$. Local propagation
of the fitted monopole uncertainty gives
$\sigma_{E_{\rm BY}/c^2}=\sigma_M/\sqrt f$ and
$\sigma_{\delta k_g}=G\sigma_M/(c^2r^2\sqrt f)$.
These are derived representations of the same fitted coefficient, not
independent measurements or a fit of the enclosing radius. If software
uses a computational default for $r$, that choice must remain explicit.
The radius-dependent boundary energy and the conserved monopole are
representations of the same spherical solution. Their contributions
are counted once. The fit does not measure a new worldtube radius, and
changing the enclosing radius in this identity is not physical mass drift.
This static local example does not assert a global cosmological embedding
or continue a static boundary observer through a horizon.

The conditional R1 extension supplies a more specific response relation
from its adopted action, frozen background and supplied matching data.
Its scalar-driven Weyl response may be added only as a distinct term
relative to the already represented reference configuration; an envelope
already absorbed into an ordinary metric multipole cannot be added again
under another name. If two parameterizations produce the same potential,
their degeneracy must be retained. In an R1 fit the reported mass is the
action's metric/source mass parameter; its relation to a point-exterior
$M_{\rm opt}$ must include the scalar contribution to the total potential.
An ordinary quadrupole is not automatically a normalized R1 boundary
amplitude. The forward transfer, state convention and nuisance freedoms
must be specified before that further interpretation is made.

Accordingly a supported mass/motion track is a reconstruction of the
effective BHSM-described optical configuration in the declared sector.
It need not distinguish that interpretation from another theory giving
the same exterior. It does not separately infer the microscopic state,
an astrophysical horizon, scalar flux selection or deformation dynamics.
Those additional inferences require their corresponding source and
matching relations; they are not conditions for retaining the ordinary
fit or for using its mass, location and motion within the stated BHSM
interpretation.
